\section{Related Work}\label{sec:related_work}

Proof of space is a mining algorithm where disk (currently SSD) space
is used instead of CPU power.
It is energetically more efficient than proof of work and
more decentralized than proof of stake.
A preliminary \emph{plotting phase}, that must be expensive, allocates large data on disk,
which is later used in a cheap
\emph{mining phase} to quickly answer
block creation challenges received from blockchain peers:
the probability of creating new blocks on top of the best chain
is proportional to the allocated data size.
Pioneering work~\cite{AtenieseBFG14,DziembowskiFKP15} introduced
formal proofs on persistent storage and showed that
on-demand recomputations (or \emph{replotting attacks}~\cite{BaigGP25}) are not viable,
due to high complexity pebbling graphs and Merkle trees:
any attempt to reduce the allocated space has a prohibitive recomputation cost.
A subsequent evolution~\cite{RenD16}
uses \emph{stacked expanders} to reduce the cost
and simplify proofs. The SpaceMint prototype~\cite{ParkKFGAP18},
now discontinued, showed the approach feasible and considered
mitigations to potential attacks, typical of
proof of space, such as mining on multiple chains and block/challenge grinding.
Chia~\cite{AbusalahACKPR17,CohenP19}, a more robust implementation, combines
space with verifiable delay functions, in a hybrid proof of space and time mechanism.
Filecoin~\cite{Fil24} lets parties commit space and prove, repeatedly,
that they still hold data in the committed space.
We build on the proof of space algorithm of Signum~\cite{signum}
based on the expensive precomputation of
plot files that later support cheap block creation, as formalized
in~\cite{Spoto25}. We work with peers built over the generic blockchain
engine Mokamint~\cite{mokamint-source-code},
derived from Signum's algorithm, but generic \wrt{} the
notion of transactions, delegated to an application layer.

There is a lack of scientific research on the use of mobile devices
for blockchain mining.
In~\cite{LeePSB20}, a system is proposed to build local
industrial IoT blockchains, where drones are mobile miners
with proof of work. It does not scale
to a global blockchain and to mainstream mobile devices, that
are not competitive against expensive specialized hardware and whose limited battery
cannot sustain the proof of work.
The same limitation applies to~\cite{SuankaewmaneeHN18}, that
presents a specialized blockchain for mobile commerce, with Android phones as
mobile miners with proof of work.
Experiments are reported for blocks holding up to six transactions, which is
sensible for a local blockchain but irrealistic for a global one.
Despite its title, \cite{JiangLW22}~tackles another problem, namely,
offloading blockchain applications
to a cloud computing provider to make them usable
with the limited resources of a phone.
More generally, this allows edge devices, with strong hardware limitations,
to offload expensive computations
to external servers and to the cloud~\cite{SoundararajS25}.
The goal of the present paper is instead to perform mining \emph{in} the mobile device
and let blockchain peers leverage many
mainstream mobile devices running a mining app.

A simple search on Google Play reports a few Android apps for
\emph{crypto mining}. Despite their name, these apps are not miners, but either
mobile dashboards to cloud pool miners (mining occurs in the cloud and
the app is just monitoring, see for instance~\cite{ctpool-google-play,viabtc-google-play,ecos-google-play,cloudminingbitcoincrypto-google-play,minex-google-play}),
or dubious games that promise Bitcoins in exchange of solving tasks
or watching ads (such as~\cite{bitcoinminerearnrealcrypto-google-play,cryptofortunetycoonbtcminer-google-play}).
In particular, our research did not find any app for actually mining \emph{in}
a mainstream mobile device.
